% =============================================================================
% ENCERRAMENTO (17:50 a 18:00)
% =============================================================================
\section[Fim]{Encerramento}

\begin{frame}{Cinco perguntas antes de confiar em uma estimativa}
  \begin{enumerate}
    \item \textbf{Qual é o contrafactual?} De onde vem e por que é crível?
    \item Os grupos comparados eram \textbf{parecidos antes}? (balanceamento, pré-tendências)
    \item \textbf{Qual pressuposto} sustenta o resultado? Foi testado ou só declarado?
    \item \textbf{Para quem vale} o efeito? (todos, tratados, quem está no corte)
    \item O que o número \textbf{não mostra}? Quem ganhou menos e por quê?
  \end{enumerate}
\end{frame}

\begin{frame}{Para levar}
  \begin{itemize}
    \item Comparar participantes e não participantes, ou antes e depois, \destaque{engana};
    \item Cada técnica \textbf{constrói um contrafactual} com um pressuposto: diga qual é e teste o que puder;
    \item A técnica certa depende da \ind{pergunta}, do \ind{desenho} da política e dos \ind{dados};
    \item Efeito médio não é tudo: \ind{avaliação qualitativa e métodos mistos} explicam como e por quê.
  \end{itemize}
\end{frame}

\begin{frame}{Para continuar: leituras (1 de 2)}
  \scriptsize\setlength{\itemsep}{0pt}
  \begin{itemize}
    \item \textbf{Batista, M. e Domingos, A. (2017).} Mais que boas intenções: técnicas quantitativas e qualitativas na avaliação de impacto de políticas públicas. \emph{Revista Brasileira de Ciências Sociais}, 32(94). \cod{doi.org/10.17666/329414/2017} \textbf{(texto-base)}
    \item Gertler, P. et al. (2011). \emph{Impact Evaluation in Practice}. World Bank.
    \item Khandker, S.; Koolwal, G.; Samad, H. (2010). \emph{Handbook on Impact Evaluation}. World Bank.
    \item Angrist, J.; Pischke, J.-S. (2009). \emph{Mostly Harmless Econometrics}. Princeton UP. \textbf{(DiD, RDD, IV)}
    \item Cunningham, S. (2021). \emph{Causal Inference: The Mixtape}. Yale UP. \textbf{(IV, controle sintético)}
    \item Abadie, A.; Diamond, A.; Hainmueller, J. (2010). Synthetic control methods for comparative case studies. \emph{JASA}, 105(490).
    \item Dunning, T. (2012). \emph{Natural Experiments in the Social Sciences}. Cambridge UP.
    \item Glennerster, R.; Takavarasha, K. (2013). \emph{Running Randomized Evaluations}. Princeton UP.
    \item Ipea / Casa Civil. \emph{Avaliação de Políticas Públicas: Guia Prático de Análise ex-Ante} e \emph{ex-Post}.
  \end{itemize}
\end{frame}

\begin{frame}{Para continuar: leituras (2 de 2)}
  \scriptsize
  \begin{itemize}
    \item Patton, M. (2002). \emph{Qualitative Research and Evaluation Methods}. Sage.
    \item Bauer, M.; Gaskell, G. (2010). \emph{Pesquisa qualitativa com texto, imagem e som}. Vozes.
    \item Braun, V.; Clarke, V. (2006). Using thematic analysis in psychology. \emph{Qualitative Research in Psychology}, 3(2).
    \item Creswell, J.; Plano Clark, V. (2018). \emph{Designing and Conducting Mixed Methods Research}. Sage.
    \item Bamberger, M. (2012). Introduction to mixed methods in impact evaluation. \emph{Impact Evaluation Notes}, 3.
    \item Beach, D.; Pedersen, R. (2013). \emph{Process-Tracing Methods}. Univ.\ of Michigan Press.
    \item Adato, M. (2008). Integrating survey and ethnographic methods to evaluate conditional cash transfer programs. IFPRI.
    \item Hunter, W.; Sugiyama, N. (2014). Transforming subjects into citizens: insights from Brazil's Bolsa Família. \emph{Perspectives on Politics}, 12(4).
  \end{itemize}
\end{frame}

\begin{frame}{Para continuar: leituras (3 de 3)}
  \scriptsize
  \begin{itemize}
    \item Ferraz, C.; Finan, F. (2008). Exposing corrupt politicians: the effects of Brazil's publicly released audits on electoral outcomes. \emph{QJE}, 123(2).
    \item Rocha, R.; Soares, R. (2010). Evaluating the impact of community-based health interventions: evidence from Brazil's Family Health Program. \emph{Health Economics}, 19(S1).
    \item Callaway, B.; Sant'Anna, P. (2021). Difference-in-differences with multiple time periods. \emph{Journal of Econometrics}, 225(2).
    \item Goodman-Bacon, A. (2021). Difference-in-differences with variation in treatment timing. \emph{Journal of Econometrics}, 225(2).
    \item Robins, J.; Rotnitzky, A.; Zhao, L. (1994). Estimation of regression coefficients when some regressors are not always observed. \emph{JASA}, 89(427).
    \item Lincoln, Y.; Guba, E. (1985). \emph{Naturalistic Inquiry}. Sage.
    \item Bardin, L. (2011). \emph{Análise de conteúdo}. Edições 70.
  \end{itemize}
\end{frame}

\begin{frame}{Repositório e contato}
  \begin{itemize}
    \item \textbf{GitHub:} \cod{github.com/andersonheri/avaliacao-pp-each}
    \item Slides e scripts em R comentados;
    \item Tudo reproduzível: mudem sementes, efeitos e tamanhos de amostra.
  \end{itemize}
  \vspace{0.1cm}
  \textbf{Prof. Dr. Anderson Henrique}\\
  {\small Pós-Doutorando no CEM/USP/CEPix; pesquisador do INCT-QualiGov e pesquisador convidado no IPEA; Doutor em Ciência Política (UFPE).\\
  \cod{andersonheri@gmail.com}}
\end{frame}

\begin{frame}{Declaração de uso de inteligência artificial}
  \begin{block}{Transparência}
    Este material foi elaborado com o apoio de uma ferramenta de inteligência artificial generativa (Claude, da Anthropic). A IA ajudou na construção dos:
    \begin{itemize}
      \item \textbf{scripts} em R e das simulações;
      \item \textbf{exemplos} usados nos slides;
      \item \textbf{imagens} dos slides (gráficos gerados pelos scripts);
      \item \textbf{checagem de gramática} e de redação.
    \end{itemize}
  \end{block}
  \footnotesize O conteúdo, as escolhas didáticas e a revisão final são do autor, que se responsabiliza pelo material. Confira as referências antes de reutilizá-las.
\end{frame}

\begin{frame}
  \centering\vspace{1.2cm}
  {\Huge\bfseries\color{azul} Obrigado!}\\[0.5cm]
  {\large Perguntas, críticas e sugestões}
\end{frame}
